
\subsection{Der Weg zur Dirac-Gleichung}
\label{subsec:weg_zur_diracgleichung}

Die Grenzen der Schrödingergleichung zeigten, dass für schnelle
Elektronen eine neue Beschreibung notwendig war. Diese musste sowohl
die Regeln der Quantenmechanik als auch die spezielle
Relativitätstheorie berücksichtigen.

Der Ausgangspunkt war die relativistische Energie-Impuls-Beziehung%
\index{Energie-Impuls-Beziehung!relativistische}

\begin{equation}
	E^2
	=
	p^2c^2
	+
	m^2c^4.
	\tag{\ref{eq:relativistische_energie_impuls_kapitel07}}
\end{equation}

Für die Energie selbst folgt daraus

\begin{equation}
	E
	=
	\pm
	\sqrt{
		p^2c^2
		+
		m^2c^4
	}.
	\label{eq:relativistische_energie_quadratwurzel}
\end{equation}

Bereits an dieser Stelle tritt ein grundlegendes Problem auf: Die
Energie besitzt einen positiven und einen negativen Lösungszweig.
Außerdem enthält die Gleichung eine Quadratwurzel aus dem
Impulsoperator.

\begin{DidacticBox}[Das zentrale Problem]
	\label{box:zentrales_problem_weg_dirac}
	
	Die Schrödingergleichung verwendet die einfache nichtrelativistische
	Energie
	
	\begin{equation}
		E_{\mathrm{kin}}
		=
		\frac{p^2}{2m}.
	\end{equation}
	
	Die relativistische Energie enthält dagegen eine Quadratwurzel:
	
	\begin{equation}
		E
		=
		\sqrt{
			p^2c^2
			+
			m^2c^4
		}.
	\end{equation}
	
	Diese Quadratwurzel lässt sich nicht ohne Weiteres in eine einfache
	lokale Wellengleichung übertragen.
\end{DidacticBox}

\subsubsection*{Von Energie und Impuls zu Operatoren}

In der Quantenmechanik werden Energie und Impuls durch Operatoren%
\index{Operator!Quantenmechanik} ersetzt.

Für die Energie gilt

\begin{equation}
	E
	\longrightarrow
	\hat{E}
	=
	\mathrm{i}\hbar
	\frac{\partial}{\partial t}.
	\label{eq:energieoperator_weg_dirac}
\end{equation}

Für den Impuls gilt

\begin{equation}
	\mathbf{p}
	\longrightarrow
	\hat{\mathbf{p}}
	=
	-\mathrm{i}\hbar\nabla.
	\label{eq:impulsoperator_weg_dirac}
\end{equation}

Dieser Übergang verbindet eine klassische Energieformel mit einer
quantenmechanischen Wellengleichung.

Bei der nichtrelativistischen Energie

\begin{equation}
	E
	=
	\frac{p^2}{2m}
\end{equation}

führt diese Ersetzung unmittelbar zur Schrödingergleichung.

Bei der relativistischen Energie ist der Weg schwieriger.

\subsubsection*{Der erste relativistische Ansatz}

Ein naheliegender Versuch besteht darin, nicht die Quadratwurzel aus
Gleichung \eqref{eq:relativistische_energie_quadratwurzel}, sondern
zunächst die quadrierte Beziehung

\begin{equation}
	E^2
	=
	p^2c^2
	+
	m^2c^4
\end{equation}

zu verwenden.

Werden Energie und Impuls durch ihre Operatoren ersetzt, entsteht die
Klein-Gordon-Gleichung\index{Klein-Gordon-Gleichung}:

\begin{equation}
	\frac{1}{c^2}
	\frac{\partial^2\Psi}{\partial t^2}
	-
	\nabla^2\Psi
	+
	\frac{m^2c^2}{\hbar^2}
	\Psi
	=
	0.
	\label{eq:klein_gordon_weg_dirac_einfach}
\end{equation}

Diese Gleichung ist relativistisch. Raum und Zeit werden in ihr
wesentlich gleichartiger behandelt als in der Schrödingergleichung.

Für Teilchen mit Spin \(0\) ist die Klein-Gordon-Gleichung eine
grundlegende Gleichung der relativistischen Quantentheorie.

Das Elektron besitzt jedoch den Spin

\begin{equation}
	s
	=
	\frac{\hbar}{2}.
\end{equation}

Daher konnte die Klein-Gordon-Gleichung das Elektron nicht vollständig
beschreiben.

\begin{PhysicsBox}[Bedeutung der Klein-Gordon-Gleichung]
	\label{box:bedeutung_klein_gordon_weg_dirac}
	
	Die Klein-Gordon-Gleichung war kein erfolgloser oder falscher Ansatz.
	
	Sie beschreibt relativistische Teilchen und Felder mit Spin \(0\).
	Für das Elektron reicht sie jedoch nicht aus, weil dessen Spin nicht in
	der Gleichung enthalten ist.
\end{PhysicsBox}

\subsubsection*{Das Problem der zweiten Zeitableitung}

Die Schrödingergleichung besitzt nur eine erste Ableitung nach der Zeit:

\begin{equation}
	\mathrm{i}\hbar
	\frac{\partial\Psi}{\partial t}
	=
	\hat{H}\Psi.
\end{equation}

Ist der Zustand \(\Psi\) zu einem bestimmten Zeitpunkt bekannt, kann
seine weitere zeitliche Entwicklung daraus berechnet werden.

Die Klein-Gordon-Gleichung enthält dagegen eine zweite Ableitung nach
der Zeit:

\begin{equation}
	\frac{\partial^2\Psi}{\partial t^2}.
\end{equation}

Dadurch unterscheidet sich ihre mathematische Struktur grundlegend von
der Schrödingergleichung.

Außerdem lässt sich die zugehörige erhaltene Dichte bei einer
Einzelteilchenbeschreibung nicht immer als positive
Wahrscheinlichkeitsdichte interpretieren.

Paul Dirac\index{Dirac, Paul} suchte deshalb nach einer relativistischen
Gleichung, die wie die Schrödingergleichung nur eine erste
Zeitableitung enthält.

\begin{DidacticBox}[Diracs Forderung]
	\label{box:diracs_forderung_erste_ordnung}
	
	Dirac wollte eine Gleichung der Form
	
	\begin{equation}
		\mathrm{i}\hbar
		\frac{\partial\Psi}{\partial t}
		=
		\hat{H}_{\mathrm D}\Psi
	\end{equation}
	
	erhalten.
	
	Der neue Hamiltonoperator sollte
	
	\begin{itemize}
		\item relativistisch korrekt,
		\item linear im Impuls und
		\item mit einer positiven Wahrscheinlichkeitsdichte vereinbar
		sein.
	\end{itemize}
\end{DidacticBox}

\subsubsection*{Die Quadratwurzel wird linearisiert}

Diracs entscheidende Idee bestand darin, die relativistische Energie
nicht direkt als Quadratwurzel zu verwenden.

Stattdessen suchte er einen Ausdruck, der linear im Impuls ist:

\begin{equation}
	E
	=
	c\boldsymbol{\alpha}\cdot\mathbf{p}
	+
	\beta mc^2.
	\label{eq:linearer_energieansatz_dirac}
\end{equation}

Dabei sind \(\boldsymbol{\alpha}\) und \(\beta\) zunächst unbekannte
Größen.

Der Ausdruck

\begin{equation}
	\boldsymbol{\alpha}\cdot\mathbf{p}
	=
	\alpha_xp_x
	+
	\alpha_yp_y
	+
	\alpha_zp_z
\end{equation}

enthält die drei räumlichen Komponenten des Impulses.

Diracs Forderung lautete: Wird der lineare Ausdruck quadriert, muss
wieder die relativistische Energie-Impuls-Beziehung entstehen:

\begin{equation}
	\left(
	c\boldsymbol{\alpha}\cdot\mathbf{p}
	+
	\beta mc^2
	\right)^2
	=
	p^2c^2
	+
	m^2c^4.
	\label{eq:forderung_quadrat_dirac_ansatz}
\end{equation}

Dieser Vorgang wird als Linearisierung%
\index{Linearisierung!relativistische Energie} bezeichnet.

\begin{PhysicsBox}[Diracs Linearisierung]
	\label{box:diracs_linearisierung}
	
	Dirac ersetzte die Quadratwurzel
	
	\begin{equation}
		E
		=
		\sqrt{
			p^2c^2
			+
			m^2c^4
		}
	\end{equation}
	
	durch einen linearen Ausdruck:
	
	\begin{equation}
		E
		=
		c\boldsymbol{\alpha}\cdot\mathbf{p}
		+
		\beta mc^2.
	\end{equation}
	
	Die unbekannten Größen \(\boldsymbol{\alpha}\) und \(\beta\) mussten so
	gewählt werden, dass beim Quadrieren wieder die relativistische
	Energieformel entsteht.
\end{PhysicsBox}

\subsubsection*{Warum gewöhnliche Zahlen nicht ausreichen}

Beim Quadrieren des Dirac-Ansatzes entstehen neben den gewünschten
Termen auch gemischte Produkte.

Vereinfacht betrachtet treten beispielsweise Ausdrücke der Form

\begin{equation}
	\alpha_x\alpha_y
	+
	\alpha_y\alpha_x
\end{equation}

und

\begin{equation}
	\alpha_x\beta
	+
	\beta\alpha_x
\end{equation}

auf.

Damit diese unerwünschten gemischten Terme verschwinden, müssen die
Größen bestimmte Bedingungen erfüllen:

\begin{equation}
	\alpha_i\alpha_j
	+
	\alpha_j\alpha_i
	=
	0
	\qquad
	\text{für}
	\quad
	i\neq j
\end{equation}

und

\begin{equation}
	\alpha_i\beta
	+
	\beta\alpha_i
	=
	0.
\end{equation}

Gewöhnliche Zahlen können diese Bedingungen nicht erfüllen, da bei
ihnen die Reihenfolge der Multiplikation keine Rolle spielt:

\begin{equation}
	ab
	=
	ba.
\end{equation}

Dirac benötigte daher mathematische Objekte, bei denen die Reihenfolge
der Multiplikation entscheidend ist.

Solche Objekte sind Matrizen\index{Matrix!Dirac-Gleichung}.

\begin{DidacticBox}[Warum erscheinen Matrizen?]
	\label{box:warum_matrizen_weg_dirac}
	
	Die Matrizen wurden nicht eingeführt, um die Gleichung künstlich
	komplizierter zu machen.
	
	Sie waren mathematisch notwendig.
	
	Nur Matrizen können die Bedingungen erfüllen, durch die beim
	Quadrieren des linearen Ansatzes alle unerwünschten gemischten Terme
	verschwinden.
\end{DidacticBox}

\subsubsection*{Vier Komponenten statt einer}

Die benötigten Matrizen besitzen mindestens vier Zeilen und vier
Spalten. Deshalb kann die Wellenfunktion nicht mehr aus nur einer
einzigen komplexen Funktion bestehen.

Sie muss vier Komponenten besitzen:

\begin{equation}
	\Psi
	=
	\begin{pmatrix}
		\psi_1 \\
		\psi_2 \\
		\psi_3 \\
		\psi_4
	\end{pmatrix}.
	\label{eq:dirac_spinor_weg_dirac}
\end{equation}

Eine solche vierkomponentige Wellenfunktion wird als
Dirac-Spinor\index{Dirac-Spinor} bezeichnet.

Die vier Komponenten enthalten mehr Information als die
Wellenfunktion der Schrödingergleichung. Sie ermöglichen die
Beschreibung von

\begin{itemize}
	\item zwei möglichen Spinrichtungen und
	\item positiven und negativen Energielösungen.
\end{itemize}

\begin{PhysicsBox}[Vom Skalar zum Spinor]
	\label{box:vom_skalar_zum_spinor}
	
	In der Schrödingergleichung kann die Wellenfunktion zunächst eine
	einzige komplexe Funktion sein.
	
	In der Dirac-Theorie besteht der Zustand aus vier miteinander
	gekoppelten Komponenten.
	
	Diese vierkomponentige Struktur ist die mathematische Grundlage für den
	Spin und die Teilchen-Antiteilchen-Beschreibung.
\end{PhysicsBox}

\subsubsection*{Der Dirac-Hamiltonoperator}

Diracs linearer Hamiltonoperator%
\index{Hamiltonoperator!Dirac-Gleichung} lautet

\begin{equation}
	\hat{H}_{\mathrm D}
	=
	c\boldsymbol{\alpha}\cdot\hat{\mathbf{p}}
	+
	\beta m_{\mathrm e}c^2.
	\label{eq:dirac_hamiltonoperator_weg}
\end{equation}

Mit dem Impulsoperator

\begin{equation}
	\hat{\mathbf{p}}
	=
	-\mathrm{i}\hbar\nabla
\end{equation}

ergibt sich

\begin{equation}
	\hat{H}_{\mathrm D}
	=
	-\mathrm{i}\hbar c
	\boldsymbol{\alpha}\cdot\nabla
	+
	\beta m_{\mathrm e}c^2.
\end{equation}

Wird dieser Hamiltonoperator in die allgemeine Gleichung der
zeitlichen Entwicklung eingesetzt, erhält man

\begin{equation}
	\mathrm{i}\hbar
	\frac{\partial\Psi}{\partial t}
	=
	\left(
	-\mathrm{i}\hbar c
	\boldsymbol{\alpha}\cdot\nabla
	+
	\beta m_{\mathrm e}c^2
	\right)
	\Psi.
	\label{eq:diracgleichung_am_ende_des_weges}
\end{equation}

Damit war die gesuchte relativistische Gleichung für das Elektron
gefunden.

\begin{HistoryBox}[Diracs Ergebnis]
	\label{box:diracs_ergebnis_weg_dirac}
	
	Dirac hatte zunächst nur versucht, Quantenmechanik und spezielle
	Relativitätstheorie miteinander zu verbinden.
	
	Die daraus entstandene Gleichung enthielt jedoch mehr als erwartet:
	
	\begin{itemize}
		\item Der Spin des Elektrons ergab sich aus der mathematischen
		Struktur.
		\item Die Gleichung besaß positive und negative Energielösungen.
		\item Die negativen Lösungen führten später zur Beschreibung des
		Positrons.
	\end{itemize}
	
	Die Mathematik sagte damit neue physikalische Eigenschaften voraus.
\end{HistoryBox}

Der entscheidende Schritt bestand somit nicht darin, die
Schrödingergleichung lediglich um eine kleine Korrektur zu ergänzen.
Dirac musste die mathematische Beschreibung des Elektrons grundlegend
erweitern.

Im nächsten Abschnitt wird die daraus entstandene Dirac-Gleichung und
ihre physikalische Bedeutung genauer betrachtet.
